\section{Introduction}
\label{sec:intro}

A frozen language model can be given new knowledge without retraining by
attaching an external memory: a store of past experience retrieved into
the context window at inference time \citep{packer2023memgpt,
park2023generative, shinn2023reflexion, zhao2024expel}, or in its strong
form a dedicated memory \emph{model} while the executive LLM stays fixed
\citep{quek2026memo}. Every such system faces the same question once it
runs long enough: \emph{what gets to stay?} Memory that only grows is
eventually dominated by stale, redundant or actively wrong entries, and
the failure is not hypothetical---memory banks an LLM continuously
rewrites degrade over time, with useful entries becoming faulty as the
rewriting accumulates \citep{zhang2026faulty}.

The prevailing answer puts a \emph{model in the loop} to curate: an LLM
rates importance \citep{park2023generative}, rewrites trajectories
\citep{zhang2026faulty}, or judges what to keep. These share two costs.
They are expensive---every curation decision is a model call---and they
introduce a \emph{proxy} between the agent and the world. A judge grades
plausibility, not consequences, and a proxy that can be optimized against
can be gamed: the reward-hacking and semantic-drift failure mode
self-training systems fall into when no external criterion anchors
quality \citep{dodgson2026survival}.

\paragraph{The curator is an attack surface, and nobody has measured it.}
Whatever plays the curator---a classifier, a judge, a heuristic, a
ledger---it decides on some signal, and that signal is reachable by an
adversary. The memory-security literature defends the store at
\emph{write} (does this text look malicious?) and at \emph{retrieval}
(should this entry be injected?) \citep{mpbench2026, memsecbench2026},
and both read content. Neither asks what happens when the input to the
\emph{deciding rule} is corrupted. The consequence is not that poison
survives; it is worse. A defence that removes entries on bad evidence,
fed bad evidence, removes the defender's own working memory. The
adversary need inject nothing at all, and the more decisively a
mechanism curates the more thoroughly it destroys the store it protects.
We call this a \emph{curation-targeted attack}.

The threat has been named. \citet{lin2026selfevolving} describe an
adversary who corrupts the performance signal to control ``the long-term
composition of the agent's cognitive state without ever writing to the
memory store directly,'' and \citet{lin2026memsurvey} give forgetting
its own lifecycle phase. Both also record that it has not been
measured---``relatively underexplored'' in the first, ``architecturally
plausible but empirically unstudied'' in the second
\citep{lin2026memsurveyv1}, whose later version drops that assessment
rather than revising it. That is the gap we
close, and it is worth being plain that closing it is the contribution:
this paper does not discover that curators are attackable, it makes the
threat quantitative.

This paper operationalizes it, measures six curation mechanisms under it,
and reports which survive. That requires a curator whose deciding signal
is not a model's opinion. Our starting observation is that a large and
practically important class of agent tasks comes with a \emph{conserved,
physically measurable resource} that already settles value without anyone
grading anything: a storage-management agent frees or wastes real bytes,
a coding agent makes a test suite pass or fail, an issue-resolution agent
resolves a GitHub issue or does not. In each case the environment returns
a scalar that is \emph{conserved}---you cannot manufacture freed bytes or
passing tests by being persuasive---and so cannot be hacked by a more
eloquent memory. The resource is conserved; the channel that
\emph{reports} it is ordinary software, and \S\ref{sec:whocanlie} says
who can write to that channel without controlling the resource. Building
on the environment-mediated selection principle of
\citet{dodgson2026survival}, we make that scalar the only currency of
memory survival.

\paragraph{\textsc{darwin-memo}.}
Each memory entry carries an \emph{energy} balance. Every cycle, every
entry pays a fixed upkeep, so existence is not free. An entry earns energy
back only when a task it decided (or supported) produces a positive
measured resource delta, with credit flowing along provenance and bounded
by a $\tanh$ so no single outcome makes an entry immortal and no single
disaster instantly kills one that has been right many times. When an
entry's energy crosses zero it is buried. The consequences fall out of the
dynamics rather than from any rule about quality: trivia nobody consults
starves on upkeep; a poisoned entry that advises a harmful action is
\emph{executed by the damage it causes} as the negative delta flows back
to it; near-duplicate survivors consolidate, pooling energy. There is no
reward model, no LLM judge, and no human curation anywhere, and because
settlement is a scalar the environment already returns, it is essentially
free.

\paragraph{The result, and its boundary.}
Under attack the mechanisms separate cleanly (\S\ref{sec:adversary}), and
what makes the difference is a bounded credit buffer with earn-back: a
lie costs energy an entry can win back, where a counter spends a life it
cannot. At one corrupted measurement per cycle the ledger keeps all of
its benign capability where every strike and quarantine policy we test
keeps at most about half, and a Thompson-sampling bandit keeps its
capability only by never removing anything, so retention is not defence.
The ledger is not exempt either, and we publish its boundary rather than
stopping at the win.

\paragraph{What we claim, and what we do not.}
It would be easy, and wrong, to claim that conserved-resource selection
dominates all alternatives. It does not, and the failures are not
marginal: a one-line heuristic matches the ledger on outcomes and, with a
real model deciding adoption, revokes poison five to nine cycles
\emph{sooner} (\S\ref{sec:wef}); a strike counter and a quarantine
policy each beat it on either side of a noise threshold in a second
environment family; a Thompson-sampling retrieval bandit reconstructed
from \citet{xu2026ael} matches it under one-sided noise; and on real
software-engineering tasks it shows no learning-curve advantage over a
token-matched random control (\S\ref{sec:swebench}). We report all of
them, and the honest summary is narrow: \emph{outside an adversarial
regime, the ledger's extra machinery buys leanness and little else.} Our
second contribution is therefore the \emph{regime map}, a pre-registered
account of which curator to reach for when.

\paragraph{Contributions.}
\begin{enumerate}[leftmargin=1.2em, itemsep=1pt, topsep=2pt]
\item \textbf{An operational threat model for an attack the literature
has named but not measured.} In a \emph{curation-targeted attack} the
adversary injects no poison and instead corrupts the signal a curator
decides on, turning the defence into the instrument that starves the
defender's benign memory. We give it a resource-bounded adversary, an
accounting identity and a budget axis, so it becomes an experiment rather
than a scenario (\S\ref{sec:adversary}), and \S\ref{sec:whocanlie} says
who can reach the settlement channel without controlling the resource
behind it.
\item \textbf{The first measurement of how curation mechanisms fare under
it}---six mechanisms at five budgets over 30 seeds---establishing that a
bounded credit buffer with earn-back is what separates the survivors,
that retention without removal is not a defence, and that every
mechanism including ours has a boundary, which we locate and publish.
\item \textbf{A label-free, judge-free curation mechanism} that grounds
survival in a conserved environmental resource, with credit assignment
along provenance and a formal description (\S\ref{sec:method}).
\item \textbf{A measured retention attack on usage-based curation},
closing a second gap the literature names without measuring
\citep{lin2026memsurvey}. A \emph{query-only} adversary---one who never
writes, never settles, and crosses no channel our threat model
names---drives two of the three components of a usage-based retention
score by asking questions, and thereby owns roughly the last $13\%$ of
the store's lifetime ($16$ of $16$ cells; $0$ of $32$ under the flat
upkeep we ship). A counterfactual swapping one denominator isolates the
cause: the exploitable property is not that the signal is \emph{usage}
but that it is \emph{relative} (\S\ref{sec:potentiation}).
\item \textbf{A sweep of the baseline family's own free parameter, and
the refutation it produced} (\S\ref{sec:countersweep}). A baseline nobody
tuned is not a baseline.
\item \textbf{The regime map} (\S\ref{sec:regime}), with control arms
reconstructed from the literature---a retrieval bandit
\citep{xu2026ael}, an LLM-judge settler, and a Generative-Agents-style
salience score \citep{park2023generative}---run on the same stdlib-only
harness rather than argued about, under pre-registered hypotheses,
seed-level bootstrap intervals, exact paired permutation tests, and
committed per-seed data validated against a manifest in CI.
\end{enumerate}
